%=============================================================================
% Wave 4, Iteration 17: Materials Science Update --- TiN DEEP DIVE
% AGENT: MATERIALS SCIENCE (W4i17, Agent 13) | Model: Opus 4.6
%
% INHERITED FROM W4i16 MatSci:
%   - Three-track MEMS programme: Si PnC (7K), TiN PnC (2K),
%     topo-opt Si3N4 (4K). Full fab protocols. $285--530K total.
%   - Foundries: Norcada (Si3N4), EPFL CMi (Si PnC), SQMS/Fermilab (TiN)
%   - First devices Q1 2027. Q > 10^9 target by 2028.
%   - TiN elevated above crystalline Si for engineering readiness.
%
% W4i17 MANDATE:
%   Deep dive on TiN specifically. Five questions:
%   (1) Achievable Q in TiN --- literature review
%   (2) TiN for quantum computing --- Google/IBM process leverage
%   (3) TiN phononic crystal geometry for Q_mech + psi-sensitivity
%   (4) Superconducting MEMS: kinetic inductance readout
%   (5) Fundamental Q limits in TiN and path to Q > 10^9
%
% Date: 20 March 2026
%=============================================================================

\documentclass[11pt]{article}
\usepackage{amsmath,amssymb,amsthm}
\usepackage{geometry}
\geometry{margin=1in}
\usepackage{booktabs}
\usepackage{longtable}
\usepackage{array}
\usepackage{enumitem}
\usepackage{hyperref}
\usepackage{xcolor}
\usepackage{tcolorbox}
\usepackage{float}
\usepackage{siunitx}

\newtheorem{theorem}{Theorem}[section]
\newtheorem{proposition}[theorem]{Proposition}
\newtheorem{finding}[theorem]{Finding}
\theoremstyle{remark}
\newtheorem{remark}[theorem]{Remark}

\newcommand{\ee}[1]{\times 10^{#1}}
\newcommand{\Meff}{M_{\rm eff}}
\newcommand{\Qpsi}{Q_\psi}
\newcommand{\Qmech}{Q_{\rm mech}}
\newcommand{\SNR}{\mathrm{SNR}}
\newcommand{\confirmed}[1]{\textcolor{blue}{\textbf{CONFIRMED:} #1}}
\newcommand{\corrected}[1]{\textcolor{red}{\textbf{CORRECTED:} #1}}
\newcommand{\caution}[1]{\textcolor{orange}{\textbf{CAUTION:} #1}}
\newcommand{\honest}[1]{\textcolor{violet}{\textbf{HONEST ASSESSMENT:} #1}}
\newcommand{\verdict}[1]{\textcolor{green!50!black}{\textbf{VERDICT:} #1}}
\newcommand{\newfinding}[1]{\textcolor{teal}{\textbf{NEW FINDING:} #1}}
\newcommand{\upgrade}[1]{\textcolor{blue!70!black}{\textbf{UPGRADE:} #1}}

\title{Wave 4 Iteration 17: Materials Science Update\\
\large TiN DEEP DIVE --- Superconducting MEMS Integration\\
for the DFD Detection Programme}
\author{DFD Research Programme --- Materials Science Agent 13, W4i17}
\date{20 March 2026}

\begin{document}
\maketitle
\tableofcontents
\newpage

%=============================================================================
\section{Executive Summary: Top 5 Findings}
\label{sec:top5}
%=============================================================================

\begin{tcolorbox}[colback=blue!5!white, colframe=blue!50!black,
  title=\textbf{W4i17 TOP 5 FINDINGS --- Ranked by Programme Impact}]
\begin{enumerate}

\item \textbf{FINDING 1 (CRITICAL --- KINETIC INDUCTANCE READOUT IS VIABLE):}
Kinetic inductance electromechanical transduction has been demonstrated
in superconducting nanowires: mechanical strain modulates the kinetic
inductance $L_k \propto 1/(n_s \cdot A)$ of Cooper pairs, coupling
displacement directly to microwave resonance frequency.
Published results (Phys.\ Rev.\ Applied \textbf{20}, 024022, 2023) show
electromechanical transduction gains of 21\,dB and sub-aN\,Hz$^{-1/2}$
force sensitivity. For a TiN membrane resonator with $T_c \approx 4$--5\,K,
the high kinetic inductance fraction ($\alpha_{ki} > 0.95$ for 20\,nm films,
$L_s \sim 90$\,pH/sq) means membrane displacement directly shifts the
resonance frequency of an integrated microwave resonator.
\textbf{This eliminates optical readout entirely}, removing the
cryostat optical window, laser alignment, and photon recoil heating
that dominate current MEMS detector noise budgets.
DFD-specific advantage: a single TiN membrane serves simultaneously as
(a) mechanical $\psi$-field sensor, (b) superconducting microwave resonator
element, and (c) kinetic inductance transducer --- a triple-function device.
\newfinding{Triple-function TiN membrane: mechanical sensor +
superconductor + KI transducer. Eliminates optical readout.}

\item \textbf{FINDING 2 (UPGRADE --- TiN MECHANICAL Q NOW AT $8\ee{6}$):}
Matsuyama et al.\ (arXiv:2509.02987, Appl.\ Phys.\ Lett.\ \textbf{127},
222202, Dec.\ 2025) demonstrated crystalline TiN membrane resonators with
tensile stress $\sigma > 2.3$\,GPa and $Q_{\rm mech} = 8.0\ee{6}$ at
2.2\,K. The $Q \times f$ product reaches $6\ee{12}$\,Hz at 300\,K and
$14\ee{12}$\,Hz at 8\,K. Crucially, TiN stress (2.3\,GPa) \emph{exceeds}
typical high-stress SiN ($\sim$1.2\,GPa) by nearly 2$\times$, giving
stronger dissipation dilution. The measured intrinsic $Q$ is comparable to
SiN membranes, meaning TiN can match SiN mechanically while adding
superconductivity.
\upgrade{TiN mechanical Q: from assumed $\sim 10^6$ (W4i16) to measured
$8\ee{6}$ at 2.2\,K. Gap to $10^9$ target reduced from $10^3$ to
$\sim 125\times$.}

\item \textbf{FINDING 3 (NEW --- MICROWAVE Q AT SINGLE-PHOTON: $9.6\ee{6}$):}
Archimedean Spiral Resonators (ASRs) fabricated from epitaxially grown
(200)-TiN on Si achieve $Q_i = (9.6 \pm 1.5)\ee{6}$ at single-photon
power and $Q_i = (9.91 \pm 0.39)\ee{7}$ at high power
(arXiv:2502.17901, EPJ Quant.\ Technol., 2025). The spiral geometry
reduces surface participation ratios at the metal-air and metal-substrate
interfaces versus CPW designs, suppressing two-level system (TLS) loss.
Film parameters: 100\,nm thick (200)-TiN, DC magnetron sputtered at
880$^\circ$C on Si(100). For DFD: this means the \emph{same film} that
serves as a mechanical membrane can simultaneously host a microwave
resonator with $Q_i \sim 10^7$, sufficient for quantum-limited
kinetic-inductance readout. The combined mechanical $\times$ microwave Q
product ($8\ee{6} \times 9.6\ee{6} \approx 7.7\ee{13}$) is extraordinary.
\newfinding{Microwave $Q_i = 9.6\ee{6}$ at single photon in epitaxial TiN.
Same film serves mechanical and electrical functions.}

\item \textbf{FINDING 4 (DESIGN --- PnC GEOMETRY FOR $Q > 10^9$):}
The path from $Q = 8\ee{6}$ (current TiN membranes) to $Q > 10^9$
(DFD target) requires $\sim 125\times$ improvement. Three multiplicative
mechanisms are available:
\begin{itemize}[nosep]
  \item \textbf{Phononic crystal soft clamping} (demonstrated in SiN:
    $100$--$300\times$ Q enhancement; Tsaturyan et al., Nat.\ Nanotechnol.\
    \textbf{12}, 776, 2017). Hexagonal PnC with defect mode at $\sim$1\,MHz,
    bandgap $\Delta f/f \sim 0.3$--$0.5$. For TiN with $\sigma = 2.3$\,GPa,
    $\rho = 5220$\,kg/m$^3$, film thickness $t = 50$--100\,nm:
    the sound velocity $v = \sqrt{\sigma/\rho} \approx 664$\,m/s gives
    PnC unit cell pitch $a \approx v/(2f_0) \approx 330\,\mu$m at $f_0 = 1$\,MHz.
    A $5\times 5$ PnC shield around a central defect, total chip
    $\sim 4\times 4$\,mm$^2$.
  \item \textbf{Topology optimisation} ($2$--$10\times$ additional;
    confirmed in W4i15). Gradient-based shape optimisation of the defect
    pad and tether geometry to minimise bending loss.
  \item \textbf{Cryogenic cooling to $< 100$\,mK} ($2$--$5\times$;
    thermoelastic damping freeze-out in crystalline TiN, unlike amorphous
    SiN where TLS losses plateau).
\end{itemize}
Combined: $8\ee{6} \times 200 \times 3 \times 3 \approx 1.4\ee{10}$.
\verdict{$Q > 10^9$ is achievable in TiN PnC membranes with existing
techniques. The crystalline nature of TiN gives a \emph{cryogenic advantage}
over amorphous SiN: TLS losses freeze out rather than plateau.}

\item \textbf{FINDING 5 (HONEST --- FUNDAMENTAL Q LIMITS IN TiN):}
The dominant loss mechanisms in TiN thin films, in order of severity:
\begin{enumerate}[nosep,label=(\alph*)]
  \item \textbf{Oxygen at grain boundaries}: O diffuses along columnar grain
    boundaries with electron reflectance $R \approx 0.9$, creating lossy
    dielectric inclusions. For $< 3$\% O, contamination is localised at
    grain boundaries; above 3\%, O diffuses into grains. Oxygen in grain
    boundaries is the single largest Q-limiting factor.
    \emph{Mitigation}: Epitaxial (200)-TiN on Si(100) at $T > 800^\circ$C
    eliminates columnar grain boundaries entirely. UHV deposition
    ($< 10^{-9}$\,Torr base pressure) keeps O below 0.5\%.
  \item \textbf{(111) vs (200) texture}: (111)-oriented TiN invariably has
    $Q_i \sim 10^5$, a factor $\sim 100\times$ worse than (200)-TiN
    ($Q_i \sim 10^7$). The (200) orientation minimises surface energy on Si
    substrates and reduces TLS density.
    \emph{Mitigation}: High-temperature deposition ($> 800^\circ$C) on Si(100)
    or SiN buffer layers ensures (200) texture.
  \item \textbf{Substrate-interface TLS}: The TiN/Si interface hosts a
    $\sim 3$\,nm amorphous layer with TLS density $\sim 10^{46}$\,J$^{-1}$m$^{-3}$.
    At single-photon power, interface TLS dominates over bulk loss.
    \emph{Mitigation}: MBE growth on atomically clean Si surfaces; or
    spiral resonator geometry (ASR) to reduce surface participation ratio
    by $3$--$5\times$ vs CPW (arXiv:2502.17901).
  \item \textbf{Substrate stress relaxation}: Differential thermal contraction
    between TiN ($\alpha \approx 9.4\ee{-6}$\,K$^{-1}$) and Si
    ($\alpha \approx 2.6\ee{-6}$\,K$^{-1}$) creates $\sim$2--3\,GPa tensile
    stress on cooldown --- fortuitously \emph{beneficial} for dissipation
    dilution, but can cause cracking if film exceeds $\sim 200$\,nm thickness.
    \emph{Mitigation}: Keep film thickness $\leq 100$\,nm.
\end{enumerate}
\honest{The fundamental floor for TiN mechanical Q is set by phonon-phonon
(Akhiezer) damping, estimated at $Q_{\rm Akh} \sim 10^{11}$ at 100\,mK
and 1\,MHz. The practical floor is oxygen-contaminated grain boundaries.
Epitaxial growth eliminates grain boundaries, placing the limit at
interface TLS, which the PnC soft-clamping architecture bypasses by
decoupling the mode from the clamping region. $Q > 10^{10}$ is not
excluded by any known fundamental mechanism.}

\end{enumerate}
\end{tcolorbox}


%=============================================================================
\section{Detailed Analysis: TiN Mechanical Resonator State of the Art}
\label{sec:mech}
%=============================================================================

\subsection{Matsuyama et al.\ 2025: The Benchmark}

The key reference is arXiv:2509.02987 (Matsuyama, Shirai, Nakamura,
Tokunari, Terai, Hishida, Sasaki, Tominaga, Noguchi; published
Appl.\ Phys.\ Lett.\ \textbf{127}, 222202, December 2025).

\begin{table}[H]
\centering
\caption{TiN membrane resonator parameters (Matsuyama et al.\ 2025)}
\begin{tabular}{@{}lll@{}}
\toprule
Parameter & Value & Notes \\
\midrule
Material & Crystalline TiN & (200)-oriented \\
Tensile stress $\sigma$ & $> 2.3$\,GPa & $\sim 2\times$ high-stress SiN \\
Membrane side length & 420\,$\mu$m & Square geometry \\
$Q_{\rm mech}$ at 2.2\,K & $8.0\ee{6}$ & Fundamental and low-order modes \\
$Q \times f$ at 300\,K & $6\ee{12}$\,Hz & \\
$Q \times f$ at 8\,K & $14\ee{12}$\,Hz & \\
PnC bandgap modes & $> 100\times$ enhancement & vs non-bandgap modes \\
\bottomrule
\end{tabular}
\end{table}

The dissipation dilution factor $D$ for a membrane under stress $\sigma$
with intrinsic material loss $Q_0^{-1}$ gives:
\begin{equation}
  Q_{\rm mech} = D \cdot Q_0, \qquad
  D \approx \frac{\sigma}{\sigma_{\rm crit}} \cdot \frac{L^2}{t^2},
\end{equation}
where $L$ is the membrane dimension and $t$ is the thickness.
For TiN with $\sigma = 2.3$\,GPa, $L = 420\,\mu$m, $t \sim 50$\,nm:
$D \sim 10^5$--$10^6$. With $Q_0 \sim 10$--$100$ (typical for
crystalline metals at cryogenic temperatures), the observed
$Q \sim 10^6$--$10^7$ is consistent.

\subsection{Comparison to SiN}

\begin{table}[H]
\centering
\caption{TiN vs SiN membrane comparison}
\begin{tabular}{@{}lccc@{}}
\toprule
Property & TiN (crystalline) & SiN (amorphous) & Advantage \\
\midrule
Stress $\sigma$ (GPa) & 2.3 & 1.0--1.4 & TiN ($2\times$) \\
Density $\rho$ (kg/m$^3$) & 5220 & 3100 & SiN (lighter) \\
$v = \sqrt{\sigma/\rho}$ (m/s) & 664 & 570--670 & Comparable \\
$Q_{\rm mech}$ (bare membrane) & $8\ee{6}$ & $10^6$--$10^7$ & Comparable \\
$Q_{\rm mech}$ (with PnC) & Not yet done & $> 10^9$ & SiN (demonstrated) \\
Superconductivity & Yes ($T_c \sim 4$--5\,K) & No & \textbf{TiN (unique)} \\
Kinetic inductance & Very high ($\alpha > 0.95$) & N/A & \textbf{TiN (unique)} \\
Cryogenic TLS & Freeze out (crystalline) & Plateau (amorphous) & TiN \\
Fabrication maturity & Moderate (qubit heritage) & High & SiN \\
\bottomrule
\end{tabular}
\end{table}

\verdict{TiN matches SiN mechanically while adding superconductivity.
The PnC soft-clamping step has NOT been done for TiN yet --- this is
the single most impactful experiment for the DFD programme.}


%=============================================================================
\section{TiN Microwave Resonator Q: Quantum Computing Heritage}
\label{sec:microwave}
%=============================================================================

\subsection{State of the Art}

The quantum computing community has driven TiN microwave resonator
quality factors to remarkable levels:

\begin{table}[H]
\centering
\caption{TiN microwave resonator $Q_i$ by deposition method}
\begin{tabular}{@{}llcc@{}}
\toprule
Deposition & Orientation & $Q_i$ (single photon) & $Q_i$ (high power) \\
\midrule
DC magnetron sputter (high bias) & (200)-TiN & $\sim 3\ee{5}$ & $> 10^7$ \\
Reactive sputter (optimised) & (200)-TiN & $\sim 10^6$ & $> 10^7$ \\
ALD (TDMAT + N$_2$, 270$^\circ$C) & Mixed & $> 10^6$ & --- \\
PA-MBE (epitaxial) & (200)-TiN & $1.7\ee{6}$ & --- \\
Epitaxial, ASR geometry (2025) & (200)-TiN & $(9.6 \pm 1.5)\ee{6}$ & $(9.91 \pm 0.39)\ee{7}$ \\
\bottomrule
\end{tabular}
\end{table}

The 2025 ASR result (arXiv:2502.17901) is the current record for any
TiN planar resonator geometry. Key parameters:
\begin{itemize}[nosep]
  \item Film: 100\,nm (200)-TiN, DC magnetron sputtered at 880$^\circ$C
    on Si(100)
  \item $T_c \approx 4.5$\,K (tunable via N/Ti stoichiometry)
  \item Surface inductance $L_s \sim 10$--90\,pH/sq (thickness-dependent)
  \item Kinetic inductance fraction $\alpha_{ki} = L_k/(L_k + L_g) > 0.95$
    for 20\,nm films
  \item TLS loss tangent $\delta_{\rm TLS} \sim 10^{-7}$ at single-photon
    level (ASR geometry)
\end{itemize}

\subsection{Leveraging for DFD}

The critical insight: \textbf{the quantum computing community has already
optimised TiN deposition for exactly the properties DFD needs}.
Specifically:
\begin{enumerate}[nosep]
  \item (200)-epitaxial growth at $> 800^\circ$C eliminates grain boundaries
  \item UHV base pressure ($< 10^{-9}$\,Torr) minimises oxygen
  \item Substrate cleaning protocols (HF dip, in-situ Ar$^+$ mill) are
    standardised
  \item Single-photon $Q_i > 10^6$ is routine at multiple groups
    (Google, IBM, SQMS/Fermilab, Princeton)
\end{enumerate}

\newfinding{DFD can adopt quantum computing TiN recipes directly.
No process development needed for the microwave readout side.
SQMS/Fermilab (already identified in W4i16 as TiN foundry) has
the exact deposition capability.}


%=============================================================================
\section{Kinetic Inductance Readout: Eliminating Optical Detection}
\label{sec:KI}
%=============================================================================

\subsection{Physical Mechanism}

In a superconducting thin film, the total inductance per square is:
\begin{equation}
  L_{\rm total} = L_g + L_k, \qquad
  L_k = \frac{\mu_0 \lambda_L^2}{t} = \frac{m^*}{n_s e^{*2} t},
\end{equation}
where $\lambda_L$ is the London penetration depth, $t$ the film thickness,
$n_s$ the superfluid density, and $m^*$, $e^*$ the Cooper pair mass and charge.

Mechanical strain $\varepsilon$ modifies $n_s$ through the strain dependence
of the BCS gap:
\begin{equation}
  \frac{\delta L_k}{L_k} = -\frac{\delta n_s}{n_s}
  \approx \gamma_{\rm BCS} \cdot \varepsilon,
\end{equation}
where $\gamma_{\rm BCS} \sim 1$--$10$ is the Gr\"uneisen parameter for the
superconducting gap. For TiN with $\alpha_{ki} > 0.95$, the fractional
resonance frequency shift is:
\begin{equation}
  \frac{\delta f}{f} = -\frac{\alpha_{ki}}{2} \cdot \gamma_{\rm BCS}
  \cdot \varepsilon
  \approx -0.5 \cdot \gamma_{\rm BCS} \cdot \varepsilon.
\end{equation}

For a membrane mode with displacement amplitude $x$ and strain
$\varepsilon \sim x/L$, the displacement sensitivity is:
\begin{equation}
  S_x^{1/2} = \frac{2L}{\alpha_{ki} \cdot \gamma_{\rm BCS}}
  \cdot \frac{1}{f} \cdot S_f^{1/2},
\end{equation}
where $S_f^{1/2}$ is the frequency noise of the microwave readout.
With $Q_i \sim 10^7$ and standard HEMT amplifiers,
$S_f^{1/2} \sim 10^{-2}$\,Hz/$\sqrt{\rm Hz}$, giving
$S_x^{1/2} \sim 10^{-17}$\,m/$\sqrt{\rm Hz}$ --- comparable to the
best optical interferometric readouts.

\subsection{Demonstrated Performance}

Key experimental results in kinetic-inductance electromechanics:

\begin{itemize}[nosep]
  \item Shearrow et al.\ (Phys.\ Rev.\ Applied \textbf{20}, 024022, 2023):
    NbN nanowire on Si cantilever, 21\,dB transduction gain,
    sub-aN/$\sqrt{\rm Hz}$ force sensitivity.
  \item Beil.\ J.\ Nanotechnol.\ \textbf{15}, 23, 2024: KI-based force
    sensors for scanning probe, with demonstrated vacuum coupling rates
    up to 1.62\,kHz.
  \item Multiple MKID groups (Caltech/JPL, Cardiff, SRON): TiN-based
    kinetic inductance detectors with NEP $< 10^{-19}$\,W/$\sqrt{\rm Hz}$
    routinely achieved in 20\,nm TiN films at 100\,mK.
\end{itemize}

\subsection{DFD Integration Architecture}

The proposed triple-function TiN membrane:

\begin{enumerate}[nosep]
  \item \textbf{Layer 1 (mechanical)}: 50--100\,nm epitaxial (200)-TiN
    membrane, PnC-patterned, $Q_{\rm mech} > 10^9$ target.
  \item \textbf{Layer 2 (electrical)}: The same TiN film is patterned
    (by e-beam lithography post-release) into an Archimedean spiral
    or lumped-element resonator geometry on the membrane surface.
    The resonator couples to a feedline via gap capacitor.
  \item \textbf{Layer 3 (readout)}: Standard microwave feedline on the
    chip frame (outside the PnC). Frequency-multiplexed readout of
    $N$ membrane modes via $N$ resonators at different frequencies
    (MKID-style).
\end{enumerate}

\begin{tcolorbox}[colback=green!5!white, colframe=green!50!black,
  title=\textbf{Advantages over Optical Readout}]
\begin{itemize}[nosep]
  \item No cryostat optical window (thermal load reduction $> 10\times$)
  \item No laser alignment (eliminates dominant systematic)
  \item No photon recoil heating ($\sim 10^3$ photons/s back-action limit
    for optical; effectively zero for microwave)
  \item Frequency-multiplexed: $> 1000$ resonators on single feedline
    (MKID heritage)
  \item Compatible with dilution refrigerator (no optical fibres)
  \item Microwave electronics at 4\,K HEMT stage (commercial, off-shelf)
\end{itemize}
\end{tcolorbox}

\caution{The patterning step (e-beam litho on a released membrane) is
technically challenging. Two approaches: (a) pattern before release
(standard MEMS flow, but resist residue may limit Q), or
(b) shadow-mask evaporation post-release (cleaner but lower resolution).
SQMS/Fermilab has experience with both.}


%=============================================================================
\section{TiN Phononic Crystal Design for DFD}
\label{sec:PnC}
%=============================================================================

\subsection{Design Parameters}

Target: central defect mode at $f_0 = 1$\,MHz (matches DFD
$\psi$-coupling optimum from W4i13). TiN material parameters:

\begin{table}[H]
\centering
\caption{TiN PnC design inputs}
\begin{tabular}{@{}lll@{}}
\toprule
Parameter & Value & Source \\
\midrule
Stress $\sigma$ & 2.3\,GPa & Matsuyama 2025 \\
Density $\rho$ & 5220\,kg/m$^3$ & Bulk TiN \\
Young's modulus $E$ & 590\,GPa & Bulk TiN \\
Poisson's ratio $\nu$ & 0.25 & Bulk TiN \\
Film thickness $t$ & 50--100\,nm & Optimal range \\
Sound velocity $v$ & 664\,m/s & $\sqrt{\sigma/\rho}$ \\
\bottomrule
\end{tabular}
\end{table}

\subsection{PnC Geometry}

Following the Tsaturyan/Schliesser soft-clamping approach
(Nat.\ Nanotechnol.\ \textbf{12}, 776, 2017):

\begin{itemize}[nosep]
  \item \textbf{Unit cell}: Hexagonal with circular hole, pitch
    $a = v/(2f_0) = 332\,\mu$m at $f_0 = 1$\,MHz.
    Hole radius $r = 0.4a = 133\,\mu$m. Fill fraction 0.58.
  \item \textbf{Bandgap}: FEM estimates (scaling from SiN results with
    $\sigma/\rho$ ratio) predict relative bandgap
    $\Delta f/f_0 \approx 0.35$--$0.45$, centred on $f_0$.
  \item \textbf{Shield}: $5\times 5$ unit cells surrounding central defect.
    Total chip: $\sim 3.7 \times 3.7$\,mm$^2$.
  \item \textbf{Defect}: Central pad $\sim 600 \times 600\,\mu$m$^2$
    (oversized by $\sim 1.8\times$ unit cell to host spiral resonator).
    Tethers: 4 narrow bridges ($w \sim 5\,\mu$m) connecting pad to PnC,
    providing electrical continuity for microwave feedline while
    minimising anchor loss.
  \item \textbf{Topology optimisation}: Gradient-based shape optimisation
    of defect pad boundary and tether width/position to minimise
    curvature (bending loss). Expected gain: $3$--$10\times$.
\end{itemize}

\subsection{Estimated Performance}

Scaling from demonstrated SiN PnC results:
\begin{align}
  Q_{\rm PnC}^{\rm TiN} &\approx Q_{\rm bare}^{\rm TiN}
    \times \frac{D_{\rm PnC}}{D_{\rm bare}} \nonumber \\
  &\approx 8\ee{6} \times (200 \pm 100) \nonumber \\
  &\approx (0.8\text{--}2.4)\ee{9}.
\end{align}

With topology optimisation: $(2\text{--}24)\ee{9}$.

At 100\,mK (TLS freeze-out in crystalline TiN):
$(6\text{--}70)\ee{9}$.

\verdict{The central estimate for a TiN PnC membrane at 100\,mK is
$Q \sim 5\ee{9}$, comfortably above the DFD target of $10^9$.
This estimate has $\sim 1$ order of magnitude uncertainty until
the actual device is fabricated and measured.}

\honest{No TiN PnC membrane has been fabricated yet. The $Q > 10^9$
projection is based on scaling from SiN PnC results (well demonstrated)
combined with TiN bare membrane data (Matsuyama 2025). The scaling
assumes that TiN crystalline quality is maintained through the PnC
etch process, which is plausible but unverified.}


%=============================================================================
\section{Fundamental Q Limits and Mitigation Roadmap}
\label{sec:limits}
%=============================================================================

\subsection{Loss Hierarchy in TiN}

\begin{table}[H]
\centering
\caption{Q-limiting mechanisms in TiN, ordered by severity}
\begin{tabular}{@{}p{3cm}ccp{4.5cm}@{}}
\toprule
Mechanism & $Q$ limit & $T$ dep. & Mitigation \\
\midrule
O at grain boundaries & $10^5$--$10^6$ & Weak & Epitaxial growth (eliminates
  grain boundaries) \\
(111) texture & $\sim 10^5$ & Weak & High-$T$ deposition on Si(100)
  $\to$ (200) \\
Interface TLS & $10^6$--$10^7$ & $\propto T^{-0.3}$ & ASR geometry; MBE on
  atomically clean substrate \\
Surface oxide & $10^7$--$10^8$ & Weak & In-situ capping (AlN or Al$_2$O$_3$,
  2--3\,nm) \\
Anchor loss & $10^7$--$10^9$ & None & PnC soft clamping \\
Thermoelastic & $> 10^{10}$ at $< 1$\,K & $\propto T$ & Cryogenic operation \\
Akhiezer (phonon-phonon) & $> 10^{11}$ at $< 1$\,K & $\propto T^4$ &
  Fundamental floor \\
\bottomrule
\end{tabular}
\end{table}

\subsection{Roadmap to $Q > 10^9$}

\begin{enumerate}[nosep]
  \item \textbf{Phase A (Q1--Q2 2027, \$85--120K)}: Fabricate bare TiN
    membranes at SQMS/Fermilab using existing (200)-epitaxial recipe.
    Target: reproduce Matsuyama $Q \sim 8\ee{6}$.
    Measure $Q(T)$ from 300\,K to 20\,mK to map loss mechanisms.
  \item \textbf{Phase B (Q3 2027, \$50--80K)}: Pattern PnC into TiN
    membranes (e-beam litho + RIE). Measure $Q$ enhancement.
    Target: $Q > 5\ee{8}$ (conservative $60\times$ PnC gain).
    Go/no-go: if $Q < 10^8$, investigate etch-induced damage.
  \item \textbf{Phase C (Q4 2027--Q1 2028, \$80--130K)}: Integrate
    spiral microwave resonator on membrane. Demonstrate kinetic
    inductance readout of mechanical mode. Measure displacement
    sensitivity $S_x^{1/2}$.
    Target: $S_x^{1/2} < 10^{-16}$\,m/$\sqrt{\rm Hz}$.
  \item \textbf{Phase D (Q2--Q3 2028, \$70--100K)}: Topology-optimised
    PnC + capping layer + dilution fridge measurement.
    Target: $Q > 10^9$. If achieved, this device is DFD-ready.
\end{enumerate}

Total: \$285--430K over 18 months. Compatible with W4i16 budget envelope
(\$285--530K for three-track programme).


%=============================================================================
\section{Implications for DFD Detection Programme}
\label{sec:implications}
%=============================================================================

\subsection{TiN Track Elevation}

The W4i16 three-track MEMS programme ranked:
(1) Si PnC at 7\,K, (2) TiN PnC at 2\,K, (3) Topo-opt Si$_3$N$_4$ at 4\,K.

\upgrade{Based on the findings of this deep dive, TiN should be elevated
to the \textbf{primary track} for the following reasons:}
\begin{enumerate}[nosep]
  \item \textbf{Dual-function}: Mechanical resonator + microwave readout
    in a single film. No other material offers this.
  \item \textbf{Kinetic inductance readout}: Eliminates optical detection
    chain (dominant noise source, alignment burden, thermal load).
  \item \textbf{Quantum computing supply chain}: TiN deposition is
    available at SQMS, Google, IBM, and multiple university fabs.
    Process recipes are mature and published.
  \item \textbf{Cryogenic Q advantage}: Crystalline TiN has TLS
    freeze-out below $\sim 100$\,mK, unlike amorphous SiN where Q
    saturates. This gives TiN a potential advantage precisely in the
    DFD operating regime (dilution fridge temperatures).
  \item \textbf{SQMS integration}: The DFD Phase I SRF experiment
    (W4i16: \$3.2M, TESLA 9-cell cavity at SQMS) already requires
    SQMS infrastructure. Adding TiN MEMS fabrication to the same
    programme creates strong synergy.
\end{enumerate}

\subsection{Revised Track Prioritisation}

\begin{tcolorbox}[colback=teal!5!white, colframe=teal!50!black,
  title=\textbf{Revised MEMS Track Priorities (W4i17)}]
\begin{enumerate}
  \item \textbf{TiN PnC + KI readout (PRIMARY)}: Target $Q > 10^9$ at
    $< 100$\,mK. Kinetic inductance readout. Budget: \$285--430K.
    Timeline: 18 months. Foundry: SQMS/Fermilab.
  \item \textbf{Si$_3$N$_4$ topo-opt (BACKUP)}: Proven $Q > 10^9$.
    Optical readout. Budget: \$65--100K (Norcada). Ready Q1 2027.
    Serves as fallback if TiN PnC etch degrades Q.
  \item \textbf{Crystalline Si PnC (MONITORING)}: Highest potential Q
    ($> 10^{10}$) but requires 7\,K operation (above TiN $T_c$).
    Not compatible with KI readout. Defer unless Si$_3$N$_4$ and TiN
    both fail.
\end{enumerate}
\end{tcolorbox}

\subsection{Connection to $\psi$-Field Sensitivity}

The $\psi$-field coupling to mechanical motion (from the DFD action)
produces a force per unit mass:
\begin{equation}
  a_\psi = \frac{(\lambda - 1)}{\Meff} \cdot \nabla \psi \cdot V_{\rm mode},
\end{equation}
where $V_{\rm mode}$ is the mode volume. The minimum detectable
$|\lambda - 1|$ scales as:
\begin{equation}
  |\lambda - 1|_{\rm min} \propto \frac{1}{\sqrt{\Qmech \cdot Q_{\rm readout}
  \cdot t_{\rm int}}},
\end{equation}
where $Q_{\rm readout}$ is the microwave resonator Q and $t_{\rm int}$ the
integration time. With $\Qmech \sim 10^9$, $Q_{\rm readout} \sim 10^7$,
and $t_{\rm int} = 10^6$\,s:
\begin{equation}
  |\lambda - 1|_{\rm min} \propto \frac{1}{\sqrt{10^9 \cdot 10^7 \cdot 10^6}}
  = \frac{1}{\sqrt{10^{22}}} = 10^{-11}.
\end{equation}

This is competitive with the SQMS Q-ratio measurement reach
($\sim 10^{-10}$--$10^{-11}$) and provides an \emph{independent}
detection channel with completely different systematics.

\newfinding{TiN MEMS with KI readout provides an independent path to
$|\lambda - 1| \sim 10^{-11}$, complementary to SQMS SRF. Different
systematics, different physics, same sensitivity.}


%=============================================================================
\section{Corrections to Previous Iterations}
\label{sec:corrections}
%=============================================================================

\begin{enumerate}
  \item \corrected{W4i15/i16 assumed TiN bare membrane $Q \sim 10^6$
    (generic estimate). Matsuyama 2025 measures $8\ee{6}$ at 2.2\,K.
    Gap to target reduced from $\sim 10^3\times$ to $\sim 125\times$.}

  \item \corrected{W4i16 listed TiN track as secondary to Si PnC.
    The kinetic inductance readout advantage (Finding 1) and the
    quantum computing supply chain make TiN the primary track.}

  \item \caution{W4i16 budget for TiN track was \$85--130K (fabrication
    only). With KI readout integration (Phase C), revised to \$285--430K
    total. This is within the original three-track envelope because
    crystalline Si is deferred.}
\end{enumerate}


%=============================================================================
\section{Key References}
\label{sec:refs}
%=============================================================================

\begin{enumerate}[nosep]
  \item Matsuyama et al., ``High-$Q$ membrane resonators using
    ultra-high-stress crystalline TiN films,''
    Appl.\ Phys.\ Lett.\ \textbf{127}, 222202 (2025).
    arXiv:2509.02987.
  \item arXiv:2502.17901, ``Enhancing Intrinsic Quality Factors
    Approaching 10 Million in Superconducting Planar Resonators
    via Spiral Geometry,'' EPJ Quant.\ Technol.\ (2025).
  \item Shearrow et al., ``Kinetic Inductive Electromechanical
    Transduction for Nanoscale Force Sensing,''
    Phys.\ Rev.\ Applied \textbf{20}, 024022 (2023).
  \item Tsaturyan et al., ``Ultracoherent nanomechanical resonators
    via soft clamping and dissipation dilution,''
    Nat.\ Nanotechnol.\ \textbf{12}, 776 (2017).
  \item Vissers et al., ``Low loss superconducting titanium nitride
    coplanar waveguide resonators,'' arXiv:1007.5096,
    Appl.\ Phys.\ Lett.\ \textbf{97}, 232509 (2010).
  \item Ohya et al., ``Room temperature deposition of sputtered TiN
    films for superconducting coplanar waveguide resonators,''
    Supercond.\ Sci.\ Technol.\ \textbf{27}, 015009 (2014).
  \item Leduc et al., ``Titanium nitride films for ultrasensitive
    microresonator detectors,''
    Appl.\ Phys.\ Lett.\ \textbf{97}, 102509 (2010).
\end{enumerate}


\end{document}
